% Self-contained report fragment. Uses amsmath/amssymb/amsthm/booktabs.
% No title, preamble, or external file dependencies.
\section{The eleven-entry cache-tested command protocol}
\label{sec:ax11qh5-proof}

This section proves the mathematical command interface obtained by combining
the eleven-entry AX splice with the coefficient-five, premise-exchanged
calculus and the cache-only Shift test. The composition needs its own proof:
the earlier eleven-entry machine tested $t+h=1$, whereas this controller tests
$h=1$. In particular, the final test of a compiled refutation must change.

The logical dependency is the inherited ordinary mathematical calibration of
the coefficient-five, premise-exchanged numerical calculus to PA. The result
below supplies explicit command identities, an all-store proof compiler,
four-field scheduling, and soundness of the changed protocol. An end-to-end
Lean or Rocq theorem is not asserted.

\subsection{Constructors and the exact AX effect}

Let
\[
 p(a,b)=\frac{(a+b)(a+b+1)}2+b,
 \qquad I(a,b)=5p(a,b),
\]
and put $M(a,b)=I(a,b)+1$, $U(i,B)=I(B,i)+2$, and
$E(a,b)=I(a,b)+3$. Write $A\to B$ for $I(A,B)$,
$N(A)=I(A,1)$, $A\wedge B=N(I(A,N(B)))$, and
$\exists_i B=N(U(i,N(B)))$. Zero is the designated truth code and one the
designated contradiction. The numerical alias $I(0,0)=0$ is retained.

The ordinary heads are
\begin{align*}
 h_1(D,H)&=(\exists_D U(D,H))\to H,\\
 h_2(D,A,H)&=U(D,U(A,H))\to U(A,U(D,H)),\\
 h_3(D,A,H)&=M(H,D)\to(E(H,A)\to M(A,D)),\\
 h_4(D,A,H)&=E(H,D)\to(M(A,H)\to M(A,D)),\\
 h_5(D,H)&=\exists_D E(H,D),\\
 h_6^{Q}(D,A,H)&=U(D,A\to H)\to(U(D,A)\to U(D,H)).
\end{align*}
The other heads are
\[
 \rho(D,A,H,u)=((H\to u)\to D)\to((D\to H)\to(A\to H)),
 \quad u\le D,
 \qquad \Gamma(B)=B\to U(B,B).
\]
Define
\begin{align*}
 V(a,b,c)&=M(a,b)\to U(c,M(a,b)),\\
 P(a,b,c)&=c\to M(c,b),\\
 L(a,b,c)&=U\bigl(b,(M(a,b)\to P(a,b,c))
                         \wedge(M(c,b)\to c)\bigr),\\
 Q(a,b,c)&=U(c,L(a,b,c)\to N(P(a,b,c))),\\
 G(a,b,c)&=N(V(a,b,c)\to Q(a,b,c)).
\end{align*}
The gate $G(a,b,c)$ is supplied only if $c\ne a,b$; otherwise its emitted
value is zero. The rules are modus ponens and generalization of theorems,
$B\mapsto U(i,B)$. Call this calculus $\mathcal D_5$.

A command store is $(t,w,d,h)$, consisting of the visible top, a Cantor
stack, an offset, and the cached formula. We write $[a;w]=p(a,w)$ and
use total projections $\operatorname{hd}0=\operatorname{tl}0=0$.
In particular, $[0;0]=0$; the notation is not a free-list encoding.

For $D=d+x$, $H=h+y$, and $A=c$, define
\[
\begin{array}{c|l}
 j&q_j\\\hline
 0&G(H,D,A)\text{ if its guard succeeds; zero otherwise}\\
 1&\Gamma(A)\\
 2&h_6^Q(D,A,H)\\
 3&h_5(D,H)\\
 4&h_4(D,A,H)\\
 5&h_3(D,A,H)\\
 6&h_2(D,A,H)\\
 7&h_1(D,H)\\
 8&\rho(D,A,H,d)\\
 9&U(d,h).
\end{array}
\]
All $q_0,\ldots,q_8$ have residue zero modulo five. The final generated
entry $q_9$ has residue two and is at least two. None of these entries is
one. Their parameters need not denote valid tagged formulas.

The spliced command has four fields:
\begin{equation}
 \mathsf A_{c,x,y,z}(t,w,d,h)
       =(0,[q_0,q_1,\ldots,q_9,t;w],z,0).
 \label{eq:ax11qh5-A}
\end{equation}
The extra push saves the guarded gate and clears the visible top. The
following full MP reader consumes $z$ from the continuation. Its input has
$t=d=h=0$, so its comparison is $I(0,z)=0$. For $z=0$ it succeeds with
conclusion zero; for $z>0$ it fails. Both branches have the effect shown
in~\eqref{eq:ax11qh5-A}. Thus soundness requires all four fields, although
the proof compiler chooses $z=0$.

\subsection{Exact source splice}

For the initial constructed association, the new source has 3,732 primitive
sites, compared with 3,714 for QH5. Old AX-return site 3613, the instruction
$\operatorname{jump}(3702)$, is replaced by a nineteen-site push at sites
3613 through 3631. Its exit 3632 is precisely the relocated full MP reader,
formerly site 3614. The inserted block agrees instruction by instruction
with the already used push at old sites 50 through 68, under an affine
relocation including its exit. All 3,713 instructions outside the replaced
site, with every successor target, agree under
\[
 r(i)=\begin{cases}i,&i\le3613,\\i+18,&i>3613.\end{cases}
\]
No external edge enters a new push interior. All fixed source nodes agree
with QH5, including its no-push MP and cache-only Shift nodes. These finite
graph facts are independently checked in
\texttt{check\_ax11qh5\_splice.py}. Reassociation of the AX grammar may
change source addresses; preservation of the complete ordinal instruction
graph is the appropriate separate check for such a change.

\subsection{All-store stack identities}

The remaining commands are, with $b=d+x$,
\[
 \mathsf S(t,w,d,h)=
 \begin{cases}\mathsf{halt},&h=1,\\
 (h,\operatorname{tl}w,d,\operatorname{hd}w),&h\ne1,
 \end{cases}
\]
and
\[
 \mathsf{MP}_x(t,w,d,h)=
 \begin{cases}(b,w,b,0),&I(t,b)=h,\\
 (0,w,b,0),&I(t,b)\ne h.
 \end{cases}
\]
Words below execute from left to right. Put $A_0=\mathsf A_{0,0,0,0}$.
Repeatedly applying the displayed Shift equation, without halting, gives
\begin{align*}
 \mathsf A_{c,x,y,0};\mathsf S^{10}(t,w,d,h)
     &=(q_8,[t;w],0,U(d,h)),\\
 \mathsf A_{c,x,y,0};\mathsf S^{11}(t,w,d,h)
     &=(U(d,h),w,0,t).
\end{align*}
Every tested cache is zero or one of the non-one generated entries.
An additional Shift now tests exactly the saved top $t$.

Define
\[
 \mathsf{Push}=A_0;\mathsf S^{10};\mathsf{MP}_0,
 \quad \mathsf{Pop}=A_0;\mathsf S^{12},
 \quad \mathsf{Discard}=\mathsf{MP}_0,
 \quad \mathsf{Offset}_i=\mathsf{MP}_i.
\]
Their full contracts are
\begin{align*}
 \mathsf{Push}(t,w,d,h)&=(0,[t;w],0,0),\\
 \mathsf{Pop}(t,[a;w],d,h)&=
   \begin{cases}\mathsf{halt},&t=1,\\(t,w,0,a),&t\ne1,\end{cases}\\
 \mathsf{Discard}(t,w,0,h)&=(0,w,0,0),\\
 \mathsf{Offset}_i(0,w,0,0)&=(0,w,i,0).
\end{align*}
Push is total even at $t=1$: its MP comparison has implication residue
zero and universal residue two, so it fails. Discard is independent of
its comparison result because both branches return zero when the target
is zero. Offset uses $I(0,i)>0$ if $i>0$, and the zero alias if $i=0$.
The Pop contract, including empty stacks and collapsed zero pushes, follows
from total projections. It differs essentially from the sum-tested AX11
Pop, which did not halt on a visible one.

For $0\le j\le8$ the selector
\[
 \mathsf{Discard};\mathsf A_{A,D,H,0};
       \mathsf S^{j+2};\mathsf{Pop}^{9-j}
 \label{eq:ax11qh5-selector}
\]
sends $(t,w,0,0)$ to $(q_j,w,0,0)$ without halting. After the initial
Shifts, the selected $q_j$ is visible; all Pops preserve it, and it is
never one. The number $9-j$ removes the remaining generated entries and
the saved zero, restoring the protected old tail. This covers the gate
and diagonal as well as the six ordinary heads.

For an actual restricted leaf with $u\le D$, use
\[
 \mathsf{Discard};\mathsf{Offset}_u;
 \mathsf A_{A,D-u,H,0};\mathsf S^{10};\mathsf{Pop}.
\]
The emitted restricted head has parameters $(D,A,H,u)$ and its valid
numerical side condition. The fourth AX field clears the offset after
emission. Generalization uses
\[
 \mathsf{Gen}_i=\mathsf{Push};\mathsf{Offset}_i;\mathsf S;
          \mathsf A_{i,0,0,0};\mathsf S^{11}.
\]
It sends $(a,w,0,0)$ to $(U(i,a),w,0,0)$ for every $a,w,i$, including
$a=1$. The intermediate Shift tests zero; AX consumes $a$ as a cache
parameter before any subsequent Shift could test it.

\subsection{Proof compilation and the final halt}

From $(0,[a,I(a,b),0;w],0,0)$ the suffix
\[
                  \mathsf S;\mathsf S;\mathsf{MP}_b;\mathsf{Pop}
\]
halts if $a=1$ or $b=1$, and otherwise returns $(b,w,0,0)$. The second
Shift tests $a$; if it continues, MP succeeds, putting $b$ in the visible
top; the Pop then tests $b$. The reserved zero below the premises is
necessary because MP itself no longer pushes a dummy cell. The same
equations cover $a=b=w=0$.

Compile a $\mathcal D_5$ derivation $e$ recursively into $W_e$. Axioms use
the selectors just proved; generalization appends $\mathsf{Gen}_i$. For an
MP node with major proof $e_1$ of $I(a,b)$ and minor proof $e_0$ of $a$, use
\[
 \mathsf{Discard};\mathsf{Push};W_{e_1};\mathsf{Push};
 W_{e_0};\mathsf{Push};\mathsf S;\mathsf S;
                   \mathsf{MP}_b;\mathsf{Pop}.
\]
Induction on the finite proof tree proves: for every $t,w$, the word
$W_e$ from $(t,w,0,0)$ either halts or returns
$(\operatorname{concl}(e),w,0,0)$. The initial Discard and Push reserve
the zero required by the MP suffix, and both recursive calls preserve the
entire protected old stack. Earlier halts satisfy the induction statement.

If $e$ proves one, the word
\[
                         A_0;W_e;\mathsf{Pop}
\]
halts from every incoming store. The first command clears the two caches,
so the induction statement applies to its resulting arbitrary stack. If
no earlier halt occurs, the final Pop tests the derived visible one.
A final plain Shift would be incorrect:
\[
               \mathsf S(1,0,0,0)=(0,0,0,0).
\]
The standalone checker \texttt{ax11qh5\_commands.py} replays 39 symbolic
identities, including the one branches, zero-stack aliases, four-field
inverse, generalization at contradiction, and the reserved-zero MP suffix.
Its symbolic values are arbitrary; any required non-one branch assumption
is explicit. It does not import the machine generator or physical simulator.

\subsection{Enumeration and soundness}

Define $e(\epsilon)=0$ and $e(c::w)=1+b_c(e(w))$, with
\[
\begin{aligned}
 b_{\mathsf S}(v)&=p(0,v),\\
 b_{\mathsf{MP}_x}(v)&=p(1,p(x,v)),\\
 b_{\mathsf A_{c,x,y,z}}(v)&=p(c+2,p(y,p(x,p(z,v)))).
\end{aligned}
\]
A fresh candidate decodes its number directly; a positive continuation
$\ell$ first decodes $\ell-1$. Consequently $b_C(e(W))$ is the fresh-candidate
code for $C$ followed by $W$, whereas $e(C::W)=1+b_C(e(W))$ is its
continuation code.

The fields are read in the physical order $y,x,z$, after the command tag
$c+2$. Total Cantor projections prove each decoding identity, including
all zero aliases. Using the old three-field AX encoding would consume the
head of the next continuation as the fourth operand; the symbolic checker
explicitly detects this error.

Every positive continuation decreases after a decoded command: the next
continuation is a component of the number beneath its initial successor.
The source arithmetic loops terminate with primitive-recursive bounds.
Consequently every nonhalting candidate finishes in finite time. The
unchanged fresh-candidate scheduler starts at two and increases its counter.
If $v=e(A_0;W_e;\mathsf{Pop})$, prefixing the halting suffix by
$\mathsf A_{k,0,0,0}$ gives the increasing unbounded family of fresh codes
\[
                  p(k+2,p(0,p(0,p(0,v)))).
\]
Each such prefix returns an ordinary store; the suffix halts from every
store. Finite candidate induction therefore proves that a refutation
forces a scheduler halt, unless the machine has halted sooner.

For the converse, annotate $t,h$ and every canonical Cantor-stack entry
with a $\mathcal D_5$ derivation. Zero has the fixed proof
$\rho(0,0,0,0)=0$; virtual entries past a zero tail use that proof.
More explicitly, $\operatorname{tl}(w)<w$ for every $w>0$.
Thus the canonical depth $\delta(w)$ is the least $j\le w$ such that
$\operatorname{tl}^{j}(w)=0$, and only the entries at depths
$0,\ldots,\delta(w)-1$ require stored annotations. If $[t;w]>0$,
its head is $t$ and its tail is $w$, so its annotation is the proof of
$t$ followed by the old canonical annotation. If $[t;w]=0$, then
$t=w=0$; there is no extra live cell, and future projections use the
fixed proof of zero. No proof or command identity counts a collapsed
zero push as a live list entry.
AX pushes only derived entries: its gate or zero, diagonal and ordinary
axioms, a restricted head satisfying $d\le d+x$, a generalization of the
old cache, and the saved old top. Its post-push MP leaves zero visible and
cached for every $z$, while the offset has no derivability requirement.
Successful MP combines two annotated premises; failure yields zero.
Shift transfers annotated entries and halts only at an annotated cache
equal to one. Thus every finite command halt yields a calculus refutation.

All word constructions, proof-tree inductions, canonical-stack decodings,
and finite candidate inductions are primitive recursive or have explicit
PA-provable correctness arguments. Composing them with the inherited
calibration gives
\[
 \mathrm{PA}\vdash
 H_{\mathrm{new\ source}}
 \leftrightarrow
 \exists p\,\operatorname{Prf}_{\mathrm{PA}}(p,\ulcorner0=1\urcorner).
\]
After the separately verified blank startup and positive physical source
simulation are attached, finite concatenation gives the same equivalence
for the exact binary table. In the converse physical direction, $T+1$
positive nonhalting source segments would extend a purported first physical
halt at time $T$, contradicting determinism. This uses finite arithmetic
induction, not a consistency assumption or a truth-reflection principle.

The coefficient-five calibration and premise-exchange theorem are inherited
mathematics. The newly checked command composition supplies the changed
eleven-entry interface and its correct final halt test. These dependency
boundaries must remain explicit in any claimed PA-provability theorem.

As a separate regression, \texttt{adversarial\_stack\_audit.py} evaluates
13,236 integer assertions using actual Cantor numbers, including both
projection identities, strict tail descent, collapsed zero pushes,
visible and cached ones, Push, Pop, generalization, and the MP suffix.
For the larger words its AX heads are small representatives of the proved
residue contracts; it is a numerical check of those command identities,
not a second check of the exact axiom emitter or the physical table.
